% Short answers to Chapter 5's exercises, lectures/L05/appendix/b_exercises.md.

\answerchapter{c5:ch}

\answer{c5:ex:1}{Tracing the signal accumulator}
\pt{a}\begin{tabular}[t]{@{}l@{\quad}c@{\quad}*{8}{c}@{}}
  \code{i} & \code{<= '0'} & 7 & 6 & 5 & 4 & 3 & 2 & 1 & 0 \\
  \midrule
  \code{bits(i)} & -- & 1 & 0 & 1 & 1 & 0 & 1 & 1 & 0 \\
  read & -- & 0 & 0 & 0 & 0 & 0 & 0 & 0 & 0 \\
  scheduled & 0 & 1 & 0 & 1 & 1 & 0 & 1 & 1 & 0 \\
  survives & no & no & no & no & no & no & no & no & yes \\
\end{tabular}
\pt{b}\code{parity_s} ends at \code{'0'}. The correct parity is \code{'1'}: \code{"10110110"} has
five \code{1}s. Every read in the pass returns the same frozen pre-pass value, and each assignment
replaces the one scheduled before it, so no iteration ever sees another's result and only the last,
\code{old_parity xor bits(0)}, survives.
\pt{c}\begin{tabular}[t]{@{}l@{\quad}*{8}{c}@{}}
  \code{i} & 7 & 6 & 5 & 4 & 3 & 2 & 1 & 0 \\
  \midrule
  \code{bits(i)} & 1 & 0 & 1 & 1 & 0 & 1 & 1 & 0 \\
  \code{acc} before & 0 & 1 & 1 & 0 & 1 & 1 & 0 & 1 \\
  \code{acc} after & 1 & 1 & 0 & 1 & 1 & 0 & 1 & 1 \\
\end{tabular}\par\smallskip
\noindent Each iteration reads the value the previous one wrote, and \code{parity} ends at the
correct \code{'1'}.

\answer{c5:ex:3}{Counting set bits}
\ptnote{Signal accumulator}The testbench stops at that input, \code{expected 1, got 2}. Every read
of the signal in the pass returns \code{1}, the count left by the previous input \code{00000001}; the
last assignment that runs, at \code{i = 1}, schedules \code{1 + 1}, and it is the one that survives.
\ptnote{Adders and cycles}Eight adders, one per unrolled iteration, chained (and then packed into
LUTs), and no clock cycles at all: the process is combinational and the count settles within one
propagation delay.
\ptnote{The range}It tells synthesis the value needs only four bits, and it is a promise the
simulator checks: if the count ever left \code{0 to 8}, GHDL would stop with a bound-check failure
at the offending line, where a \code{std_logic_vector(3 downto 0)} would hold \code{9} to \code{15}
without a word.

\answer{c5:ex:4}{A synchronizer module}
\pt{b}\begin{itemize}
  \item That every button is being held down, until the real released level arrives.
  \item The line falling low straight after reset: a start bit, so the receiver would start
    receiving a frame that was never sent.
  \item \code{'0'} is a flip-flop's ordinary reset value, as in every other reset in the book, but no
    default can know what a signal means: both real users override it, so getting \code{PRESET}
    right is the instantiating design's job, and it should always pass it explicitly.
\end{itemize}
\pt{c}Eight flip-flops at \code{SIZE = 4}, two at \code{SIZE = 1}. \code{PRESET} ends up as each
flip-flop's reset value, whether its asynchronous reset clears or presets it. Both generics were
used up before synthesis placed anything: \code{SIZE} decided how many flip-flops to build and how
wide the ports are, \code{PRESET} how each one resets.

\answer{c5:ex:5}{What Quartus does to your VHDL}
\pt{a}\begin{itemize}
  \item It is the root of the design: its ports become the chip's pins, and every other entity is
    built only where it is instantiated beneath it.
  \item The testbench has no ports to put on pins, and its body (timed waits, assertions, reports)
    describes a simulation, not hardware, so it cannot be synthesized; only \code{min_of_two} can.
\end{itemize}
\pt{b}\begin{itemize}
  \item The pins are an input to fitting: placement, routing and the programming file are all made
    for them, and a \file{.sof} compiled before a change is not patched afterwards.
  \item To whatever pins the fitter picked on its own, not the LEDs you meant; some project settings
    refuse to fit unassigned pins at all.
\end{itemize}
\pt{c}\begin{itemize}
  \item The same circuit for both, one four-input LUT holding the truth table of
    \code{x = ab + c'd}: the internal signal only names a wire, and synthesis builds the function,
    not the text.
  \item At two terms, honestly either; the named version pulls ahead as soon as the expression
    outgrows one line, and the habit is worth having before you need it.
\end{itemize}
\pt{d}\begin{itemize}
  \item A clean C build does what the code says; legal VHDL can still describe hardware you did not
    mean, such as a latch, a combinational loop, an unassigned pin or an incomplete sensitivity list,
    and Quartus reports all of those as warnings.
  \item An inferred latch.
\end{itemize}

\answer{c5:ex:6}{Predicting parity}
\pt{a}\code{"00000000"}: \code{'0'}; \code{"10110110"}: \code{'1'}; \code{"11111111"}: \code{'0'};
\code{"10000000"}: \code{'1'}.
\pt{b}It sweeps all 256 inputs, these four among them, and prints only
\code{parity_gen: all checks passed!}.
\pt{c}\begin{itemize}
  \item It inverts.
  \item Parity is the XOR of all eight bits, and XORing in a flipped bit complements the result:
    the number of \code{1}s changes by exactly one, so odd becomes even or even odd.
  \item An even number of flipped bits, two for instance: each inverts the parity, so the second
    undoes the first and the parity still matches.
\end{itemize}

\answer{c5:ex:7}{Reviewing a latched decoder}
\pt{a}A latch on each of the four bits of \code{leds}. The \code{when others} branch assigns
nothing, so on that path VHDL requires \code{leds} to keep its old value; remembering is something
combinational logic cannot do, so a latch is the only thing the tool can build.
\pt{b}\begin{itemize}
  \item \code{sel} is a \code{std_logic_vector}, each bit with nine values, so \code{when others}
    matches 78 of the 81 values it can hold (\code{'U'} before it is driven, \code{'X'}, \code{'Z'}
    and the rest), not just \code{"11"}.
  \item The latch is inferred from the code, not from the values that occur: the process still has
    a path that leaves \code{leds} unassigned, so the tool builds the latch whatever \code{sel}
    does, and the design contains storage nobody asked for.
\end{itemize}
\pt{c}In the branch, \code{when others => leds <= "0000";}. Before the \code{case}, the default
\code{leds <= "0000";}. The second is easier to maintain: the default covers every path, including
branches added later, where the first has to be got right again each time the \code{case} grows.
Both give the same combinational logic.
\pt{d}An error stops you; a warning lets the design through, and because the latch usually holds the
right value the design appears to work, so the bug ships and surfaces later as an intermittent
fault in a design that ``compiled cleanly''.
